\chapter{Refinamento de Consultas e Views}
\label{cap:refinamento_consultas_views}

À medida que o domínio da SQL avança, o estudante percebe que não basta apenas selecionar colunas, aplicar filtros e juntar tabelas. Em aplicações reais, torna-se necessário refinar resultados, controlar ordenação, limitar a quantidade de linhas devolvidas, eliminar duplicidades, combinar consultas distintas, produzir análises por grupo sem perder a granularidade das linhas e encapsular consultas complexas para reutilização.

Esta aula aprofunda exatamente esse nível de sofisticação. O foco recai sobre cláusulas e recursos avançados do \texttt{SELECT}, bem como sobre o conceito de \texttt{VIEW}, que permite transformar consultas recorrentes em objetos lógicos reutilizáveis dentro do banco de dados.

\section{Por que refinar consultas}

No uso cotidiano de bancos de dados, surgem demandas como:

\begin{itemize}
    \item mostrar registros em ordem específica;
    \item paginar listas extensas;
    \item exibir apenas valores distintos;
    \item combinar resultados oriundos de fontes diferentes;
    \item comparar cada linha com o desempenho médio de seu grupo;
    \item simplificar consultas repetidas em relatórios ou sistemas.
\end{itemize}

Tais necessidades fazem com que o comando \texttt{SELECT} deixe de ser apenas uma ferramenta de recuperação simples e passe a atuar como mecanismo de organização, filtragem refinada, síntese e apresentação dos dados.

\section{ORDER BY}

A cláusula \texttt{ORDER BY} define a ordem em que as linhas serão apresentadas. Sem ela, não se deve assumir que os resultados surgirão sempre na mesma sequência.

\subsection{Ordenação simples}

\begin{sqlcode}
SELECT nome, idade
FROM alunos
ORDER BY nome;
\end{sqlcode}

Nesse caso, os alunos serão exibidos em ordem alfabética pelo nome.

\subsection{Ordenação decrescente}

\begin{sqlcode}
SELECT nome, idade
FROM alunos
ORDER BY idade DESC;
\end{sqlcode}

Agora, os alunos são listados da maior para a menor idade.

\subsection{Ordenação por múltiplos critérios}

\begin{sqlcode}
SELECT nome, curso, idade
FROM alunos
ORDER BY curso ASC, idade DESC;
\end{sqlcode}

Essa consulta ordena primeiro pelo curso e, dentro de cada curso, coloca os alunos mais velhos antes dos mais jovens.

\section{LIMIT}

A cláusula \texttt{LIMIT} restringe a quantidade de linhas retornadas pela consulta.

\begin{sqlcode}
SELECT *
FROM alunos
ORDER BY aluno_id
LIMIT 5;
\end{sqlcode}

Esse comando é útil para:
\begin{itemize}
    \item exibir apenas uma amostra;
    \item reduzir o volume de saída em testes;
    \item implementar paginação;
    \item limitar relatórios resumidos.
\end{itemize}

\section{OFFSET}

A cláusula \texttt{OFFSET} determina quantas linhas iniciais serão ignoradas antes de retornar o conjunto desejado.

\begin{sqlcode}
SELECT *
FROM alunos
ORDER BY aluno_id
LIMIT 5 OFFSET 5;
\end{sqlcode}

Nesse caso, o banco ignora as primeiras cinco linhas e devolve as cinco seguintes.

\subsection{Paginação}

\texttt{LIMIT} e \texttt{OFFSET} são frequentemente usados juntos para paginação.

Exemplo:
\begin{itemize}
    \item página 1: \texttt{LIMIT 10 OFFSET 0}
    \item página 2: \texttt{LIMIT 10 OFFSET 10}
    \item página 3: \texttt{LIMIT 10 OFFSET 20}
\end{itemize}

\begin{quadro_dica}
Sempre que usar \texttt{LIMIT} e \texttt{OFFSET}, use também \texttt{ORDER BY}. Caso contrário, a noção de “primeiras” ou “próximas” linhas perde sentido lógico.
\end{quadro_dica}

\section{DISTINCT}

A cláusula \texttt{DISTINCT} elimina duplicidades do resultado.

\begin{sqlcode}
SELECT DISTINCT curso
FROM alunos;
\end{sqlcode}

Se cem alunos estiverem distribuídos em apenas três cursos, o resultado trará apenas três linhas, uma para cada curso distinto.

\subsection{DISTINCT em mais de uma coluna}

\begin{sqlcode}
SELECT DISTINCT curso, status
FROM alunos;
\end{sqlcode}

Nesse caso, a eliminação de duplicidades ocorre sobre a combinação dos valores de \texttt{curso} e \texttt{status}.

\section{UNION}

O operador \texttt{UNION} combina os resultados de duas consultas compatíveis, eliminando duplicidades entre elas.

\begin{sqlcode}
SELECT nome
FROM alunos
UNION
SELECT nome
FROM professores;
\end{sqlcode}

Para que a operação seja válida:
\begin{itemize}
    \item ambas as consultas devem devolver o mesmo número de colunas;
    \item as colunas correspondentes devem possuir tipos compatíveis.
\end{itemize}

\section{UNION ALL}

O \texttt{UNION ALL} também combina resultados, mas preserva todas as linhas, inclusive as repetidas.

\begin{sqlcode}
SELECT nome, 'ALUNO' AS origem
FROM alunos
UNION ALL
SELECT nome, 'PROFESSOR' AS origem
FROM professores;
\end{sqlcode}

Esse operador é especialmente útil quando:
\begin{itemize}
    \item as duplicidades têm significado;
    \item deseja-se melhor desempenho;
    \item é importante saber a origem de cada linha.
\end{itemize}

\section{HAVING revisitado}

Embora já estudado anteriormente, o \texttt{HAVING} se conecta naturalmente ao refinamento das consultas.

\begin{sqlcode}
SELECT curso_id, COUNT(*) AS total
FROM alunos
GROUP BY curso_id
HAVING COUNT(*) >= 3
ORDER BY total DESC;
\end{sqlcode}

Aqui, primeiro os dados são agrupados, depois os grupos são filtrados e, por fim, ordenados.

\section{Funções de janela}

As funções de janela (\textit{window functions}) permitem calcular valores analíticos sobre um conjunto de linhas relacionadas à linha atual, mas sem colapsar o resultado como ocorre com \texttt{GROUP BY}.

\subsection{Média por grupo em cada linha}

\begin{sqlcode}
SELECT
    aluno_id,
    nome,
    curso_id,
    idade,
    AVG(idade) OVER (PARTITION BY curso_id) AS media_idade_curso
FROM alunos;
\end{sqlcode}

Cada linha continua representando um aluno, mas agora traz ao lado a média de idade do seu curso.

\subsection{Classificação por grupo}

\begin{sqlcode}
SELECT
    nome,
    curso_id,
    idade,
    RANK() OVER (PARTITION BY curso_id ORDER BY idade DESC) AS posicao
FROM alunos;
\end{sqlcode}

Essa consulta classifica os alunos dentro de cada curso, do mais velho para o mais novo.

\subsection{Contagem por grupo sem perder linhas}

\begin{sqlcode}
SELECT
    nome,
    curso_id,
    COUNT(*) OVER (PARTITION BY curso_id) AS total_no_curso
FROM alunos;
\end{sqlcode}

\section{Diferença entre GROUP BY e funções de janela}

A distinção é central:

\begin{itemize}
    \item \texttt{GROUP BY} resume várias linhas em uma linha por grupo;
    \item funções de janela preservam cada linha individual e acrescentam informações analíticas.
\end{itemize}

Essa diferença torna as funções de janela especialmente úteis para relatórios em que se deseja manter o detalhe sem abrir mão da comparação com o grupo.

\section{Views}

\subsection{O que é uma view}

Uma \texttt{VIEW} é uma tabela virtual baseada em uma consulta. Ela não precisa armazenar os dados como uma tabela comum; em vez disso, armazena a lógica que produz aquele resultado.

\subsection{Por que usar views}

As views ajudam a:
\begin{itemize}
    \item simplificar consultas repetitivas;
    \item encapsular junções complexas;
    \item facilitar o uso por outros usuários;
    \item servir de base para relatórios;
    \item contribuir com estratégias de segurança.
\end{itemize}

\subsection{Criando uma view simples}

\begin{sqlcode}
CREATE VIEW vw_alunos_cursos AS
SELECT
    a.nome AS aluno,
    c.nome AS curso
FROM alunos a
INNER JOIN cursos c
    ON a.curso_id = c.curso_id;
\end{sqlcode}

\subsection{Consultando a view}

\begin{sqlcode}
SELECT *
FROM vw_alunos_cursos
ORDER BY aluno;
\end{sqlcode}

\section{Views como abstração}

Uma das principais vantagens das views é abstrair a complexidade estrutural. Em vez de obrigar o usuário a lembrar e reescrever junções repetidas, a lógica já fica encapsulada.

\begin{observacao}
Uma view é especialmente útil quando a mesma consulta passa a ser utilizada repetidamente em exercícios, relatórios, painéis ou rotinas da aplicação.
\end{observacao}

\section{Views e segurança}

As views também podem ser usadas para restringir o que será exposto a determinados usuários. Em vez de conceder acesso direto à tabela original, pode-se conceder acesso apenas à view, que mostra um recorte controlado das informações.

Por exemplo, uma view pode ocultar:
\begin{itemize}
    \item colunas sensíveis;
    \item dados administrativos;
    \item atributos desnecessários a determinados perfis.
\end{itemize}

\section{Exemplo resolvido 1 -- Consulta ordenada e paginada}

Suponha que a secretaria queira exibir uma lista de alunos em ordem alfabética, mostrando apenas os dez primeiros de cada página.

\begin{sqlcode}
SELECT nome, curso_id
FROM alunos
ORDER BY nome
LIMIT 10 OFFSET 0;
\end{sqlcode}

\textbf{Interpretação:}
\begin{itemize}
    \item os resultados são organizados alfabeticamente;
    \item apenas dez linhas são mostradas;
    \item essa lógica pode ser usada na primeira página de um relatório.
\end{itemize}

\section{Exemplo resolvido 2 -- Lista de cursos distintos}

Suponha que se deseje apenas saber quais cursos aparecem entre os alunos cadastrados.

\begin{sqlcode}
SELECT DISTINCT curso_id
FROM alunos
ORDER BY curso_id;
\end{sqlcode}

\textbf{Interpretação:}
\begin{itemize}
    \item duplicidades são eliminadas;
    \item o resultado é uma lista enxuta;
    \item esse tipo de consulta ajuda em filtros e menus.
\end{itemize}

\section{Exemplo resolvido 3 -- União de perfis}

Imagine um relatório institucional que deva exibir nomes de alunos e professores em uma única listagem.

\begin{sqlcode}
SELECT nome, 'ALUNO' AS origem
FROM alunos
UNION ALL
SELECT nome, 'PROFESSOR' AS origem
FROM professores;
\end{sqlcode}

\textbf{Interpretação:}
\begin{itemize}
    \item as duas consultas possuem a mesma estrutura;
    \item o campo \texttt{origem} permite distinguir os papéis;
    \item \texttt{UNION ALL} preserva todas as linhas.
\end{itemize}

\section{Exemplo resolvido 4 -- Média do grupo em cada linha}

Deseja-se listar alunos com a média de idade do curso ao lado.

\begin{sqlcode}
SELECT
    nome,
    curso_id,
    idade,
    AVG(idade) OVER (PARTITION BY curso_id) AS media_curso
FROM alunos;
\end{sqlcode}

\textbf{Interpretação:}
\begin{itemize}
    \item cada linha mantém um aluno individual;
    \item a média do curso é calculada sem reduzir o resultado a uma linha por curso;
    \item isso facilita comparações detalhadas.
\end{itemize}

\section{Exemplo resolvido 5 -- View para relatório acadêmico}

\begin{sqlcode}
CREATE VIEW vw_notas_alunos AS
SELECT
    a.nome AS aluno,
    d.nome AS disciplina,
    n.nota
FROM notas n
INNER JOIN alunos a
    ON n.aluno_id = a.aluno_id
INNER JOIN disciplinas d
    ON n.disciplina_id = d.disciplina_id;
\end{sqlcode}

\begin{sqlcode}
SELECT *
FROM vw_notas_alunos
ORDER BY aluno, disciplina;
\end{sqlcode}

\textbf{Interpretação:}
\begin{itemize}
    \item a consulta mais complexa é centralizada em uma view;
    \item usuários podem trabalhar diretamente com a visão consolidada;
    \item a leitura do relatório se torna mais amigável.
\end{itemize}

\section{Erro comum 1 -- Usar LIMIT sem ORDER BY}

Exemplo problemático:

\begin{sqlcode}
SELECT *
FROM alunos
LIMIT 10;
\end{sqlcode}

Sem ordenação explícita, o conjunto de “dez primeiros” pode não fazer sentido lógico do ponto de vista do relatório.

\section{Erro comum 2 -- Confundir DISTINCT com GROUP BY}

\texttt{DISTINCT} elimina duplicidades. \texttt{GROUP BY} organiza linhas em grupos para permitir agregações. Eles não são equivalentes.

\section{Erro comum 3 -- Usar UNION quando UNION ALL seria mais adequado}

\texttt{UNION} elimina duplicidades e pode custar mais. Se as repetições não forem problema, \texttt{UNION ALL} pode ser a melhor escolha.

\section{Erro comum 4 -- Tentar substituir toda agregação por função de janela}

Funções de janela não substituem todos os usos de \texttt{GROUP BY}. Elas têm papel complementar, não absoluto.

\section{Erro comum 5 -- Criar view sem objetivo claro}

Uma view deve ter propósito claro: simplificação, segurança, relatório ou reutilização. Criar views sem critério pode apenas aumentar a complexidade do esquema.

\section{Boas práticas}

\begin{itemize}
    \item usar \texttt{ORDER BY} sempre que a ordem importar;
    \item combinar \texttt{LIMIT} com ordenação;
    \item usar \texttt{DISTINCT} apenas quando a eliminação de duplicidades fizer sentido;
    \item escolher conscientemente entre \texttt{UNION} e \texttt{UNION ALL};
    \item nomear views de forma clara e padronizada;
    \item empregar funções de janela quando se deseja preservar detalhe com contexto analítico.
\end{itemize}

\begin{quadro_dica}
Uma convenção útil é prefixar views com \texttt{vw\_}. Isso ajuda a distingui-las rapidamente de tabelas comuns.
\end{quadro_dica}

\section{Minicaso integrador -- Painel acadêmico resumido}

Considere que a coordenação deseja montar um painel com as seguintes informações:

\begin{enumerate}
    \item lista paginada de alunos;
    \item cursos distintos existentes;
    \item união entre nomes de alunos e professores;
    \item média de idade do curso ao lado de cada aluno;
    \item uma view para relatório de notas.
\end{enumerate}

Uma sequência possível seria:

\subsection*{Passo 1 -- Lista paginada}

\begin{sqlcode}
SELECT nome, curso_id
FROM alunos
ORDER BY nome
LIMIT 10 OFFSET 0;
\end{sqlcode}

\subsection*{Passo 2 -- Cursos distintos}

\begin{sqlcode}
SELECT DISTINCT curso_id
FROM alunos
ORDER BY curso_id;
\end{sqlcode}

\subsection*{Passo 3 -- União de nomes}

\begin{sqlcode}
SELECT nome, 'ALUNO' AS origem
FROM alunos
UNION ALL
SELECT nome, 'PROFESSOR' AS origem
FROM professores;
\end{sqlcode}

\subsection*{Passo 4 -- Média por curso em cada linha}

\begin{sqlcode}
SELECT
    nome,
    curso_id,
    idade,
    AVG(idade) OVER (PARTITION BY curso_id) AS media_curso
FROM alunos;
\end{sqlcode}

\subsection*{Passo 5 -- View de relatório}

\begin{sqlcode}
CREATE VIEW vw_notas_alunos AS
SELECT
    a.nome AS aluno,
    d.nome AS disciplina,
    n.nota
FROM notas n
INNER JOIN alunos a
    ON n.aluno_id = a.aluno_id
INNER JOIN disciplinas d
    ON n.disciplina_id = d.disciplina_id;
\end{sqlcode}

\begin{sqlcode}
SELECT *
FROM vw_notas_alunos
ORDER BY aluno, disciplina;
\end{sqlcode}

\textbf{Comentário didático:}
Esse minicaso mostra como vários recursos de refinamento podem coexistir em um mesmo cenário. A consulta deixa de ser apenas mecanismo de busca e passa a atuar como ferramenta de apresentação, organização e reutilização.

\section{Exercícios básicos}

\begin{enumerate}
    \item Escreva uma consulta ordenada por nome.
    \item Escreva uma consulta ordenada por idade em ordem decrescente.
    \item Faça uma consulta que retorne apenas cinco alunos.
    \item Liste os cursos distintos.
    \item Escreva uma consulta que una nomes de alunos e professores.
\end{enumerate}

\section{Exercícios intermediários}

\begin{enumerate}
    \item Monte uma paginação com \texttt{LIMIT} e \texttt{OFFSET}.
    \item Compare \texttt{UNION} e \texttt{UNION ALL} com um exemplo.
    \item Crie uma consulta com \texttt{HAVING} e \texttt{ORDER BY}.
    \item Escreva uma função de janela para calcular média por curso.
    \item Crie uma view que simplifique a junção entre alunos e cursos.
\end{enumerate}

\section{Exercícios avançados}

\begin{enumerate}
    \item Classifique alunos por idade dentro de cada curso usando função de janela.
    \item Crie duas views distintas para relatórios acadêmicos.
    \item Explique em que contexto uma view pode ajudar na segurança.
    \item Monte uma consulta paginada com filtro, ordenação e limitação de resultados.
\end{enumerate}

\section{Exercício integrador}

Considere um sistema universitário com as tabelas \texttt{alunos}, \texttt{cursos}, \texttt{professores}, \texttt{disciplinas} e \texttt{notas}. Elabore:

\begin{enumerate}
    \item uma consulta que liste cursos distintos;
    \item uma consulta ordenada por curso e nome;
    \item uma consulta paginada de alunos;
    \item uma consulta que una nomes de alunos e professores;
    \item uma função de janela que calcule média por curso;
    \item uma view para uso em relatório acadêmico;
    \item uma breve justificativa para o uso de cada recurso.
\end{enumerate}

\section{Síntese da aula}

Esta aula aprofundou o refinamento do \texttt{SELECT} ao estudar \texttt{ORDER BY}, \texttt{LIMIT}, \texttt{OFFSET}, \texttt{DISTINCT}, \texttt{UNION}, \texttt{UNION ALL}, funções de janela e \texttt{VIEW}. Com exemplos resolvidos, erros comuns, exercícios progressivos e um minicaso integrador, o estudante amplia sua capacidade de produzir consultas mais úteis, legíveis e reutilizáveis. A próxima aula tratará do controle de acesso, roles e permissões.